% page 425 du fac-simile (image p-424 du scan)
% bloc 165.1 x 233.3 mm ; marge gauche 36.9 mm
\pgc{15.240mm}{0.000mm}{
\l{\cel{49}417}
\l{}
\l{\sou{palestro}. (en\cel{1}\sou{la epoki antiqu}a).Exerci korpala (lukto, boxo,}
\l{\cel{3}salto, kuro, e c.) - DEFIRS.}
\l{}
\l{\sou{paleto}.\cel{2}Parto plata di la remilo, qua frapas la aquo ; o di}
\l{\cel{3}roto navala, aquo-falala. - F.}
\l{}
\l{\sou{palio}. Stipito di planti cereala di qua onu deprenis la grani.}
\l{\cel{3}- FIS.}
\l{}
\l{\sou{paliatar}. (\sou{trans.})\cel{2}Mingravigar (morbo) sen remediar lu komplete.}
\l{\cel{3}DEFIRS.}
\l{}
\l{\sou{palimpsesto}. (\sou{paleogr.)}\cel{3}Manuskripto di qua la dat-unesma}
\l{\cel{3}skriburo gratesis, skrapesis, por posibligar la skribo di}
\l{\cel{3}nova texto (ulakaze : transverse). - DEFIRS.}
\l{}
\l{\sou{palinodio}. (en\cel{1}\sou{la epoki antiqua}).\cel{2}Poemo en qua onu retraktas}
\l{\cel{3}to quon onu skribis en poemo antea. -\cel{2}DEFIS.}
\l{}
\l{\sou{paliso}. Sproso ligna, di qua onu igis pintoza un extremajo,}
\l{\cel{3}e de qua onu facas (ordinare) palisado. - EFIS.}
\l{}
\l{\sou{palisado}. I. Klozilo ek planki, stangi, e c. - lineo de ligno-}
\l{\cel{3}peci triangula destinita ad obstruktar la fauco di konstrukt-}
\l{\cel{3}uro militala - ad impedar la avanco di la enemiko ye la bordo}
\l{\cel{3}ti fosato. - II. (\sou{gardeno)}. Lineo de arbori di qui onu lasas}
\l{\cel{3}kreskar la branchi de la pedo, tale ke li formacas quaza hego.}
\l{\cel{3}- DEFIS.}
\l{}
\l{\sou{palisandro}. Ligno violea, kun nuanci flava, nigra, quan onu}
\l{\cel{3}povas polisar adbrile, uzata por la inkrusturi, por facar}
\l{\cel{3}mobli. - L.\cel{1}\sou{dalbergia nigra},\cel{1}\sou{jacarania mimosifoli}a. DEFIRS.}
\l{}
\l{\sou{paliumo}. I. (en\cel{1}\sou{la epoki antiqu}a). Nomo per qua la Romani signi-}
\l{\cel{3}fikis la mantelo di la Greki.-II. (\sou{eklezio katol}.) Bendo}
\l{\cel{3}de lano blanka, ornita ye quar kruci, qua cirkondas la shul-}
\l{\cel{3}tri e pendas, avane e dope, quan la papo +weras o quan lu}
\l{\cel{3}sendas a la arkiepiskopi, primati, e c. kom signo pri judicio-}
\l{\cel{3}yuro. - DEFIS.}
\l{}
\l{\sou{palmo}. I. Brancho di palmiero, di qua la folii dispozesis quale}
\l{\cel{3}manuo apertata. - II. La interna parto di la manuo. -DEFIRS.}
\l{}
\l{\sou{palmata.}\cel{2}\sou{(zool.)}\cel{3}Di qua la fingri interligesas per membrano.}
\l{\cel{3}- EFIS.}
\l{}
\l{\sou{palmeto}. Mikra palmo, ornivo skultita, sur frisi, e c. -}
\l{}
\l{\sou{palmipedo}. (\sou{zool.}) Klaso de uceli di qui la pedi esas palmata}
\l{\cel{3}(anado, ansero, e c.) - EFIS.}
\l{}
\l{palpar. (trans.)\cel{2}Explorar per la manuo, quale la blindi. - EFIS.}
}
